% ==============================================================================
% Plantilla Institucional de Trabajo Terminal — UPIIZ IPN
% Archivo: ejemplos/tablas.tex
% ==============================================================================
% Ejemplos prácticos para tablas simples, autoajustables (tabularx) y multipágina (longtable).
% ==============================================================================

\section{Guía de ejemplos para tablas}
\label{sec:ejemplo-tablas}

% ------------------------------------------------------------------------------
% 1. TABLA PEQUEÑA / ESTÁNDAR CON BOOKTABS
% ------------------------------------------------------------------------------
En la \tabref{tab:ejemplo-simple} se ilustra una tabla estándar con líneas limpias mediante el paquete \texttt{booktabs}.

\begin{table}[htbp]
    \centering
    \caption{Ejemplo de tabla compacta con formato formal.}
    \label{tab:ejemplo-simple}
    \begin{tabular}{llr}
        \toprule
        \textbf{Componente} & \textbf{Modelo} & \textbf{Consumo (\unit{\watt})} \\
        \midrule
        Microcontrolador & ESP32-WROOM-32 & 0.8 \\
        Controlador H-Bridge & L298N Dual & 25.0 \\
        Sensor IMU & MPU-6050 & 0.02 \\
        Servomotor & MG996R & 4.5 \\
        \bottomrule
    \end{tabular}
\end{table}

% ------------------------------------------------------------------------------
% 2. TABLA ANCHA AUTOAJUSTABLE AL ANCHO DE PÁGINA (tabularx)
% ------------------------------------------------------------------------------
Para tablas con texto descriptivo largo que deban ocupar exactamente el ancho de página, utilice la columna tipo \texttt{X} de \texttt{tabularx} (\tabref{tab:ejemplo-tabularx}).

\begin{table}[htbp]
    \centering
    \caption{Ejemplo de tabla ajustable con columnas de justificación automática.}
    \label{tab:ejemplo-tabularx}
    \begin{tabularx}{\textwidth}{l c c X}
        \toprule
        \textbf{Subsistema} & \textbf{Voltaje (\unit{\volt})} & \textbf{Corriente (\unit{\ampere})} & \textbf{Función y Criterio de Selección} \\
        \midrule
        Alimentación & 12.0 & 5.0 & Batería LiPo de 3 celdas seleccionada por alta densidad energética y bajo peso. \\
        Regulación & 5.0 & 3.0 & Convertidor Buck DC-DC con eficiencia superior al $90\,\%$ para evitar pérdidas térmicas. \\
        Sensado & 3.3 & 0.2 & Sensores digitales I2C protegidos con conversores de nivel lógico bidireccionales. \\
        \bottomrule
    \end{tabularx}
\end{table}
